\documentclass[10pt,a4paper]{amsart}
\usepackage{amsmath,amssymb,amsthm,graphicx}
% Retained preamble for the Stacks Project chapter morphisms.tex (Morphisms of Schemes).
% The chapter file pinned-original-source.tex carries only \input{preamble} on its first line, so no
% preamble text lives inside the chapter itself.  The Stacks Project uses one single shared preamble for
% every chapter; the snapshot stages it verbatim in sources/stacks-chow/preamble.tex.  The block below is
% copied mechanically, line for line, from that file's lines 14-18 (the Xy-pic setup) and lines 154-217
% (the native theorem declarations together with the mathematical macros).  Nothing was added, removed or
% redefined.  The wrapper class and the amsmath/amssymb/amsthm/graphicx base packages are supplied by the
% assembler, so the class/\externaldocument part of the shared preamble is not repeated here.
% For commutative diagrams we use Xy-pic
\usepackage[all]{xy}

% We use 2cell for 2-commutative diagrams.
\xyoption{2cell}
% Theorem environments.
%
\theoremstyle{plain}
\newtheorem{theorem}[subsection]{Theorem}
\newtheorem{proposition}[subsection]{Proposition}
\newtheorem{lemma}[subsection]{Lemma}

\theoremstyle{definition}
\newtheorem{definition}[subsection]{Definition}
\newtheorem{example}[subsection]{Example}
\newtheorem{exercise}[subsection]{Exercise}
\newtheorem{situation}[subsection]{Situation}

\theoremstyle{remark}
\newtheorem{remark}[subsection]{Remark}
\newtheorem{remarks}[subsection]{Remarks}

\numberwithin{equation}{subsection}

% Macros
%
\def\lim{\mathop{\mathrm{lim}}\nolimits}
\def\colim{\mathop{\mathrm{colim}}\nolimits}
\def\Spec{\mathop{\mathrm{Spec}}}
\def\Hom{\mathop{\mathrm{Hom}}\nolimits}
\def\Ext{\mathop{\mathrm{Ext}}\nolimits}
\def\SheafHom{\mathop{\mathcal{H}\!\mathit{om}}\nolimits}
\def\SheafExt{\mathop{\mathcal{E}\!\mathit{xt}}\nolimits}
\def\Sch{\mathit{Sch}}
\def\Mor{\mathop{\mathrm{Mor}}\nolimits}
\def\Ob{\mathop{\mathrm{Ob}}\nolimits}
\def\Sh{\mathop{\mathit{Sh}}\nolimits}
\def\NL{\mathop{N\!L}\nolimits}
\def\CH{\mathop{\mathrm{CH}}\nolimits}
\def\proetale{{pro\text{-}\acute{e}tale}}
\def\etale{{\acute{e}tale}}
\def\QCoh{\mathit{QCoh}}
\def\Ker{\mathop{\mathrm{Ker}}}
\def\Im{\mathop{\mathrm{Im}}}
\def\Coker{\mathop{\mathrm{Coker}}}
\def\Coim{\mathop{\mathrm{Coim}}}

% Boxtimes
%
\DeclareMathSymbol{\boxtimes}{\mathbin}{AMSa}{"02}

%
% Macros for moduli stacks/spaces
%
\def\QCohstack{\mathcal{QC}\!\mathit{oh}}
\def\Cohstack{\mathcal{C}\!\mathit{oh}}
\def\Spacesstack{\mathcal{S}\!\mathit{paces}}
\def\Quotfunctor{\mathrm{Quot}}
\def\Hilbfunctor{\mathrm{Hilb}}
\def\Curvesstack{\mathcal{C}\!\mathit{urves}}
\def\Polarizedstack{\mathcal{P}\!\mathit{olarized}}
\def\Complexesstack{\mathcal{C}\!\mathit{omplexes}}
% \Pic is the operator that assigns to X its picard group, usage \Pic(X)
% \Picardstack_{X/B} denotes the Picard stack of X over B
% \Picardfunctor_{X/B} denotes the Picard functor of X over B
\def\Pic{\mathop{\mathrm{Pic}}\nolimits}
\def\Picardstack{\mathcal{P}\!\mathit{ic}}
\def\Picardfunctor{\mathrm{Pic}}
\def\Deformationcategory{\mathcal{D}\!\mathit{ef}}

% --- appended 2026-10-12 (Chief): the project own definitions of the reference/slogan/history
% display environments, copied mechanically from the shared Stacks preamble lines 8-12
% (\usepackage{verbatim} plus the three \newenvironment lines, which the chapter preamble keeps
% verbatim from the project preamble). Nothing added, removed or redefined: these environments
% carry no mathematical content and the project itself typesets them as comments. Without these
% three lines a counted statement whose own text uses them fails with Environment slogan undefined.
\usepackage{verbatim}
\newenvironment{reference}{\comment}{\endcomment}
\newenvironment{slogan}{\comment}{\endcomment}
\newenvironment{history}{\comment}{\endcomment}
\usepackage{cleveref}
\crefname{originallemma}{original lemma}{original lemmas}
\Crefname{originallemma}{Original lemma}{Original lemmas}
\makeatletter
\newlabel{algebra-lemma-image-dense-generic-points}{{00FL}{source}{}{sourcepointer-1}{}}
\newlabel{algebra-lemma-image-dense-generic-points@cref}{{[originallemma][0][]00FL}{[1][0][]source}{}{}{}}
\makeatother
\crefname{originallemma}{original lemma}{original lemmas}
\Crefname{originallemma}{Original lemma}{Original lemmas}
\makeatletter
\newlabel{lemma-reduction-universal-homeomorphism}{{054M}{source}{}{sourcepointer-2}{}}
\newlabel{lemma-reduction-universal-homeomorphism@cref}{{[originallemma][0][]054M}{[1][0][]source}{}{}{}}
\makeatother
\crefname{originalproposition}{original proposition}{original propositions}
\Crefname{originalproposition}{Original proposition}{Original propositions}
\makeatletter
\newlabel{proposition-universal-homeomorphism}{{0CNE}{source}{}{sourcepointer-3}{}}
\newlabel{proposition-universal-homeomorphism@cref}{{[originalproposition][0][]0CNE}{[1][0][]source}{}{}{}}
\makeatother
\title{For a dominant ring map there is a subalgebra whose spectrum map is a universal homeomorphism and which contains the images of all such factorizations}
\author{The Stacks Project Authors}
\date{stacks-project:morphisms.tex:local-snapshot-20261012; GNU FDL 1.2 or later}
\begin{document}
\maketitle
\noindent\textit{Editorial note.} All delivered mathematical text is the author's own, copied from the pinned Stacks Project chapter file pinned-original-source.tex; no statement or proof step was rewritten, reordered or removed.\par\medskip
\noindent\textit{Editorial note.} Following the multi-claim convention for this source, source\_url is the source-level URL https://stacks.math.columbia.edu/ and the source label lemma-relative-weak-normalization-algebra is carried in claim\_key and here rather than in a per-claim tag URL.\par\medskip
\noindent\textit{Editorial note.} The preamble is the chapter's shared Stacks Project preamble, retained as sources/stacks-morphisms/retained-preamble.tex: a line-for-line excerpt of the shared Stacks preamble staged with this snapshot. No macro was invented or redefined, and this card adds no preamble text of its own.\par\medskip
\noindent\textit{Editorial note.} Statement span 14747-14758 and proof span 14760-14810 come from the halfA pool card operations/pool-cards/stacks-morphisms-halfA.json and were re-read in the pinned chapter file before this card was written.\par\medskip
\noindent\textit{Editorial note.} Tag self-check per LANE\_BRIEF 10.5 (5): the label lemma-relative-weak-normalization-algebra resolves in the project's own tag list to canonical tag 0H3J (line 21319 of the index, full label morphisms-lemma-relative-weak-normalization-algebra). reference/public-items.jsonl holds no row whose primary source\_file is morphisms.tex, so the intersection with the same-chapter public tag ids is empty and the candidate passes.\par\medskip
\noindent\textit{Editorial note.} The gate was run as: python operations/agent-authoring/tools/precard\_gate.py --source morphisms.tex --label lemma-relative-weak-normalization-algebra --start 14747 --end 14758 --proof-start 14760 --proof-end 14810 --local-id 34093, giving RESULT=OK with all four span guards passing, no public interval row for this chapter and occupancy FREE for morphisms.tex 14747 14810. Spans were recomputed live from the label in the pinned chapter file and the located text matched the pool byte counts. The proof span contains a commutative diagram drawn with xymatrix, which the retained 88-line chapter preamble supports.\par\medskip
\noindent\textit{Editorial note.} 3 cross-references occur inside the counted spans; they are listed in external\_references as navigation pointers only.\par\medskip
\noindent\textit{Allowed standard prerequisites (editorial).}\begin{itemize}
\item Reduction of universal homeomorphism statements to the associated reduced rings, cited from the same chapter. Cited at proof line 14761. Lets the whole argument be carried out after replacing both rings by their reductions.
\item For a reduced ring map inducing a dominant map of spectra the base is contained in the target, cited from the algebra chapter. Cited at proof lines 14779-14780. Makes the reduced case an inclusion of rings, which is what the structure result for universal homeomorphisms requires as input.
\item The structure of a finite type algebra inducing a universal homeomorphism: every element lies in an extension generated by finitely many elements satisfying square or cube relations or p-th power relations over the preceding ones, cited from the same chapter. Cited at proof lines 14783-14790. Shows that the image of any admissible factorization is contained in the subalgebra constructed below, which is property (2) of the statement.
\end{itemize}\medskip
\ifdefined\hypertarget\hypertarget{sourcepointer-1}{}\fi
\noindent\textit{Source reference.} Original lemma 00FL: focused-seeds/AGP/sources/stacks-snapshot-20261011/tags.txt, source lines 575--575. Same-chapter reference used in the retained author text at source line(s) 1301, 11842, 11946, 12203, 14779. It is supplied only as a navigation pointer so that the retained text compiles; the counted item delivers the author's own statement and the author's complete proof, and no part of the referenced item is copied in or substituted for the counted argument.\par\medskip
\ifdefined\hypertarget\hypertarget{sourcepointer-2}{}\fi
\noindent\textit{Source reference.} Original lemma 054M: focused-seeds/AGP/sources/stacks-snapshot-20261011/tags.txt, source lines 6645--6645. Same-chapter reference used in the retained author text at source line(s) 11413, 11943, 14761. It is supplied only as a navigation pointer so that the retained text compiles; the counted item delivers the author's own statement and the author's complete proof, and no part of the referenced item is copied in or substituted for the counted argument.\par\medskip
\ifdefined\hypertarget\hypertarget{sourcepointer-3}{}\fi
\noindent\textit{Source reference.} Original proposition 0CNE: focused-seeds/AGP/sources/stacks-snapshot-20261011/tags.txt, source lines 15889--15889. Same-chapter reference used in the retained author text at source line(s) 11950, 14781, 14808, 14874. It is supplied only as a navigation pointer so that the retained text compiles; the counted item delivers the author's own statement and the author's complete proof, and no part of the referenced item is copied in or substituted for the counted argument.\par\medskip
\begin{lemma}
\label{lemma-relative-weak-normalization-algebra}
Let $A \to B$ be a ring map inducing a dominant morphism
$\Spec(B) \to \Spec(A)$ of spectra. There exists an $A$-subalgebra
$B' \subset B$ such that
\begin{enumerate}
\item $\Spec(B') \to \Spec(A)$ is a universal homeomorphism,
\item given a factorization $A \to C \to B$ such that
$\Spec(C) \to \Spec(A)$ is a universal homeomorphism, the
image of $C \to B$ is contained in $B'$.
\end{enumerate}
\end{lemma}

\begin{proof}
We will use Lemma \ref{lemma-reduction-universal-homeomorphism}
without further mention. Consider the commutative diagram
$$
\xymatrix{
B \ar[r] & B_{red} \\
A \ar[u] \ar[r] & A_{red} \ar[u]
}
$$
For any factorization $A \to C \to B$ of $A \to B$ as in (2), we see
that $A_{red} \to C_{red} \to B_{red}$ is a factorization of
$A_{red} \to B_{red}$ as in (2). It follows that if the
lemma holds for $A_{red} \to B_{red}$ and produces the $A_{red}$-subalgebra
$B'_{red} \subset B_{red}$, then setting $B' \subset B$
equal to the inverse image of $B'_{red}$ solves the lemma for
$A \to B$. This reduces us to the case discussed in the next paragraph.

\medskip\noindent
Assume $A$ and $B$ are reduced. In this case $A \subset B$ by
Algebra, Lemma \ref{algebra-lemma-image-dense-generic-points}.
Let $A \to C \to B$ be a factorization as in (2).
Then we may apply Proposition \ref{proposition-universal-homeomorphism}
to $A \subset C$ to see that every element of $C$ is contained
in an extension $A[c_1, \ldots, c_n] \subset C$ such that for
$i = 1, \ldots, n$ we have
\begin{enumerate}
\item $c_i^2, c_i^3 \in A[c_1, \ldots, c_{i - 1}]$, or
\item there exists a prime number $p$ with
$pc_i, c_i^p \in A[c_1, \ldots, c_{i - 1}]$.
\end{enumerate}
Thus property (2) holds if we define $B' \subset B$ to be the
subset of elements $b \in B$ which are contained in an extension
$
A[b_1, \ldots, b_n] \subset B
$
such that (*) holds: for $i = 1, \ldots, n$ we have
\begin{enumerate}
\item $b_i^2, b_i^3 \in A[b_1, \ldots, b_{i - 1}]$, or
\item there exists a prime number $p$ with
$pb_i, b_i^p \in A[b_1, \ldots, b_{i - 1}]$.
\end{enumerate}
There are only two things to check: (a) $B'$ is an $A$-subalgebra, and
(b) $\Spec(B') \to \Spec(A)$ is a universal homeomorphism.
Part (a) follows because given $n \geq 0$ and $b_1, \ldots, b_n \in B$
satisfying (*) and $m \geq 0$ and $b'_1, \ldots, b'_m \in B$
satisfying (*), the integer
$n + m$ and $b_1, \ldots, b_n, b'_1, \ldots, b'_m \in B$
also satisfies (*). Finally, part (b) holds by
Proposition \ref{proposition-universal-homeomorphism} and our construction
of $B'$.
\end{proof}

\end{document}
